% GENERATED FILE. Edit canonical lessons, then run npm run generate:books.
% canonical-source-hash: fnv1a64:80530bd2130f7031
% canonical-lessons: HI-C463-sukun, HI-C463-santos, HI-C463-tasalli, HI-C463-aghana, HI-C463-santvana

\chapter{Calm Again}
\label{ch:hi-calm-again}
\hlchaptermodality{463}

\begin{chapteropening}
\textbf{By the end of this chapter:} \emph{I can say that someone felt calm and content.}

Complete the last lesson of chapter 463: I can say that someone felt calm and content.
\end{chapteropening}

\section[\texorpdfstring{\hi{सुकून} (sukūn)}{सुकून (sukūn)}]{\hi{सुकून} (sukūn) — peace of mind}

\label{lesson:HI-C463-sukun}



\begin{quote}
Before the new one: say the Hindi for to strut, then the Hindi for fussiness.
\end{quote}

\subsection*{You'll want to know: \hi{सुकून}}
\textbf{\hi{सुकून}} — \emph{sukūn} — \textquotedblleft{}peace of mind\textquotedblright{}.

\begin{quote}
\hi{गाँव} \hi{की} \hi{शांत} \hi{शाम} \hi{में} \hi{बहुत} \hi{सुकून} \hi{है।}
\end{quote}

\begin{quote}
\emph{gā̃v kī śānt śām meṁ bahut sukūn hai.}
\end{quote}

\begin{quote}
There is a great peace in a quiet village evening.
\end{quote}

Masculine.

\begin{cousinweb}
\textbf{\hi{सुकून}} (\emph{sukūn}) is usually traced to Arabic \emph{sukūn}, \textquotedblleft{}stillness, rest\textquotedblright{}, which reached Hindi through Persian.
\end{cousinweb}

\begin{grammarlens}[title={What you've built}]
One more word for describing what you lived and felt.
\end{grammarlens}

\subsection*{Your turn}
\begin{itemize}
\raggedright
  \item \emph{Say it:} \emph{sukūn}
  \item \emph{Say it:} \emph{sukūn}, once more
  \item \emph{Say it:} \emph{gā̃v kī śānt śām meṁ bahut sukūn hai.} — There is a great peace in a quiet village evening.
  \item \emph{Recall:} say the Hindi for to strut, then the Hindi for fussiness, then say \emph{sukūn} again
\end{itemize}

\begin{tcolorbox}[breakable,colback=teal!4,colframe=teal!35!black,title={Before you move on}]
What does \hi{सुकून} mean? (\textquotedblleft{}Peace of mind\textquotedblright{}.) Say it once more.
\end{tcolorbox}
\section[\texorpdfstring{\hi{संतोष} (santoṣ)}{संतोष (santoṣ)}]{\hi{संतोष} (santoṣ) — contentment; satisfaction}

\label{lesson:HI-C463-santos}



\begin{quote}
Before the new one: say the Hindi for fussiness, then the Hindi for peace of mind.
\end{quote}

\subsection*{You'll want to know: \hi{संतोष}}
\textbf{\hi{संतोष}} — \emph{santoṣ} — \textquotedblleft{}contentment; satisfaction\textquotedblright{}.

\begin{quote}
\hi{दादा} \hi{को} \hi{अपने} \hi{जीवन} \hi{से} \hi{पूरा} \hi{संतोष} \hi{है।}
\end{quote}

\begin{quote}
\emph{dādā ko apne jīvan se pūrā santoṣ hai.}
\end{quote}

\begin{quote}
Grandfather is completely content with his life.
\end{quote}

Masculine.

\begin{cousinweb}
\textbf{\hi{संतोष}} (\emph{santoṣ}) is a tatsama, borrowed from Sanskrit \emph{saṁtoṣa}, \textquotedblleft{}satisfaction, contentedness\textquotedblright{} (Monier-Williams).
\end{cousinweb}

\begin{grammarlens}[title={What you've built}]
One more word for describing what you lived and felt.
\end{grammarlens}

\subsection*{Your turn}
\begin{itemize}
\raggedright
  \item \emph{Say it:} \emph{santoṣ}
  \item \emph{Say it:} \emph{santoṣ}, once more
  \item \emph{Say it:} \emph{dādā ko apne jīvan se pūrā santoṣ hai.} — Grandfather is completely content with his life.
  \item \emph{Recall:} say the Hindi for fussiness, then the Hindi for peace of mind, then say \emph{santoṣ} again
\end{itemize}

\begin{tcolorbox}[breakable,colback=teal!4,colframe=teal!35!black,title={Before you move on}]
What does \hi{संतोष} mean? (\textquotedblleft{}Contentment; satisfaction\textquotedblright{}.) Say it once more.
\end{tcolorbox}
\section[\texorpdfstring{\hi{तसल्ली} (tasallī)}{तसल्ली (tasallī)}]{\hi{तसल्ली} (tasallī) — reassurance; comfort}

\label{lesson:HI-C463-tasalli}



\begin{quote}
Before the new one: say the Hindi for peace of mind, then the Hindi for contentment.
\end{quote}

\subsection*{You'll want to know: \hi{तसल्ली}}
\textbf{\hi{तसल्ली}} — \emph{tasallī} — \textquotedblleft{}reassurance; comfort\textquotedblright{}.

\begin{quote}
\hi{डाक्टर} \hi{की} \hi{बात} \hi{सुनकर} \hi{परिवार} \hi{को} \hi{तसल्ली} \hi{हुई।}
\end{quote}

\begin{quote}
\emph{ḍākṭar kī bāt sunkar parivār ko tasallī huī.}
\end{quote}

\begin{quote}
The family was reassured by what the doctor said.
\end{quote}

Feminine.

\begin{cousinweb}
\textbf{\hi{तसल्ली}} (\emph{tasallī}) is usually traced to Arabic \emph{tasallī}, \textquotedblleft{}consolation\textquotedblright{}, which reached Hindi through Persian.
\end{cousinweb}

\begin{grammarlens}[title={What you've built}]
One more word for describing what you lived and felt.
\end{grammarlens}

\subsection*{Your turn}
\begin{itemize}
\raggedright
  \item \emph{Say it:} \emph{tasallī}
  \item \emph{Say it:} \emph{tasallī}, once more
  \item \emph{Say it:} \emph{ḍākṭar kī bāt sunkar parivār ko tasallī huī.} — The family was reassured by what the doctor said.
  \item \emph{Recall:} say the Hindi for peace of mind, then the Hindi for contentment, then say \emph{tasallī} again
\end{itemize}

\begin{tcolorbox}[breakable,colback=teal!4,colframe=teal!35!black,title={Before you move on}]
What does \hi{तसल्ली} mean? (\textquotedblleft{}Reassurance; comfort\textquotedblright{}.) Say it once more.
\end{tcolorbox}
\section[\texorpdfstring{\hi{अघाना} (aghānā)}{अघाना (aghānā)}]{\hi{अघाना} (aghānā) — to eat one's fill; to be content}

\label{lesson:HI-C463-aghana}



\begin{quote}
Before the new one: say the Hindi for contentment, then the Hindi for reassurance.
\end{quote}

\subsection*{You'll want to know: \hi{अघाना}}
\textbf{\hi{अघाना}} — \emph{aghānā} — \textquotedblleft{}to eat one's fill; to be content\textquotedblright{}.

\begin{quote}
\hi{दादी} \hi{के} \hi{हाथ} \hi{का} \hi{खाना} \hi{खाकर} \hi{बच्चे} \hi{अघा} \hi{गए।}
\end{quote}

\begin{quote}
\emph{dādī ke hāth kā khānā khākar bacce aghā gae.}
\end{quote}

\begin{quote}
After my grandmother's home cooking, the children were full and content.
\end{quote}

Past (he …): \hi{वह} \hi{खाकर} \hi{अघा} \hi{गया}, \emph{vah khākar aghā gayā}.

\begin{cousinweb}
\textbf{\hi{अघाना}} (\emph{aghānā}) is a tadbhava: CDIAL lists Hindi \emph{aghānā}, \textquotedblleft{}to be full; satiated\textquotedblright{}, under Sanskrit \emph{āghrāpayati} (CDIAL 1062). One is \emph{aghā} when hunger, or a wish, is fully met.
\end{cousinweb}

\begin{grammarlens}[title={What you've built}]
One more word for describing what you lived and felt.
\end{grammarlens}

\subsection*{Your turn}
\begin{itemize}
\raggedright
  \item \emph{Say it:} \emph{aghānā}
  \item \emph{Say it:} \emph{aghānā}, once more
  \item \emph{Say it:} \emph{dādī ke hāth kā khānā khākar bacce aghā gae.} — After my grandmother's home cooking, the children were full and content.
  \item \emph{Recall:} say the Hindi for contentment, then the Hindi for reassurance, then say \emph{aghānā} again
\end{itemize}

\begin{tcolorbox}[breakable,colback=teal!4,colframe=teal!35!black,title={Before you move on}]
What does \hi{अघाना} mean? (\textquotedblleft{}To eat one's fill; to be content\textquotedblright{}.) Say it once more.
\end{tcolorbox}
\section[\texorpdfstring{\hi{सांत्वना} (sāntvanā)}{सांत्वना (sāntvanā)}]{\hi{सांत्वना} (sāntvanā) — consolation; words of comfort}

\label{lesson:HI-C463-santvana}



\begin{quote}
Before the new one: say the Hindi for reassurance, then the Hindi for to eat one's fill.
\end{quote}

\subsection*{You'll want to know: \hi{सांत्वना}}
\textbf{\hi{सांत्वना}} — \emph{sāntvanā} — \textquotedblleft{}consolation; words of comfort\textquotedblright{}.

\begin{quote}
\hi{दोस्तों} \hi{की} \hi{सांत्वना} \hi{से} \hi{मीरा} \hi{का} \hi{मन} \hi{शांत} \hi{हुआ।}
\end{quote}

\begin{quote}
\emph{dostoṁ kī sāntvanā se Mīrā kā man śānt huā.}
\end{quote}

\begin{quote}
Her friends' words of comfort calmed Mira's mind.
\end{quote}

Feminine.

\begin{cousinweb}
\textbf{\hi{सांत्वना}} (\emph{sāntvanā}) is a tatsama, borrowed from Sanskrit \emph{sāntvanā}, which Monier-Williams explains as calming someone with gentle words, or consolation (Monier-Williams). It is what one offers a person in grief: \emph{sāntvanā denā}.
\end{cousinweb}

\begin{grammarlens}[title={What you've built}]
One more word for describing what you lived and felt.
\end{grammarlens}

\subsection*{Your turn}
\begin{itemize}
\raggedright
  \item \emph{Say it:} \emph{sāntvanā}
  \item \emph{Say it:} \emph{sāntvanā}, once more
  \item \emph{Say it:} \emph{dostoṁ kī sāntvanā se Mīrā kā man śānt huā.} — Her friends' words of comfort calmed Mira's mind.
  \item \emph{Recall:} say the Hindi for reassurance, then the Hindi for to eat one's fill, then say \emph{sāntvanā} again
\end{itemize}

\begin{tcolorbox}[breakable,colback=teal!4,colframe=teal!35!black,title={Before you move on}]
What does \hi{सांत्वना} mean? (\textquotedblleft{}Consolation; words of comfort\textquotedblright{}.) Say it once more.
\end{tcolorbox}
